%قبل از شروع مسئله ترجیحاً چیزی ننویسید. اما اگر خواستید بنویسید، در قالب لیستی (فقرات و شمارش) بنویسید.
%برای اضافه کردن نمره، بعدی از شروع‌مسئله، [x] را اضافه کنید و به جای x  نمره مورد نظر را قرار دهید. برای زیرمسئله‌ها هم به همین صورت است.
%\شروع‌مسئله[10]

\شروع‌مسئله[]
حداکثر 8تا

زیرگراف شامل اعداد ۱۰۰۱ تا ۱۰۲۴ را در نظر بگیرید، چون که هرم یک هرم کمینه است، و ۲۴ عدد بزرگترین اعداد این هرم می‌باشند، زیرگراف حاصل همبند و شامل ریشه می‌باشد پس خود یک درخت مانند درخت کمینه بدون شرط نقریبا کامل بودن می‌باشد. این گراف را $T$ می‌نامیم.

لم: اگر در ارتفاع $h$ از یک درخت دودویی $k$ رأس باشد در ارتفاع $h-1$ حداقل
$\ceil{\frac{k}{2}}$
رأس می‌باشد.

اثبات: فرض کنید که لم برقرار نباشد. هر رأس  از ارتفاع $h$ یک پدر از ارتفاع $h-1$ دارد، لذا اگر رأس‌های این دسته را بر حسب پدر‌هایشان دسته‌بندی کنیم، طبق اصل لانه‌ی کبوتری حداقل یک دسته با بیش از سه عضو وجود دارد که مخالف فرض دودویی بودن درخت می‌باشد.

چون که
$1024 = 2^{10}$
با کنار گذاشتن یک رأس درخت تبدیل به یک درخت کامل با ارتفاع 10 می‌شود. پس همه‌ی برگ‌های درخت غیر از یکی ارتفاع ۱۰ دارند.

حال درخت $T$ را در نظر بگیرید، هدف بیشینه کردن تعداد برگ‌های این درخت در عمق 10 می‌باشد.

ادعا می‌کنیم که این تعداد کمتر‌ از 9 واحد می‌باشد.
 
 اگر این تعداد حداقل 9 واحد باشد، در ارتفاع‌های ۱ تا ۱۰ به ترتیب حداقل
$[1, 1, 1, 1, 1, 1, 2, 3, 5, 9]$
رأس باید موجود باشد. یعنی حداقل 25 رأس.

برای 8 رأس مثال زیر از بین پاسخ‌ها انتخاب شده‌است. (با تشکر از محمدصدرا حیدری)

\begin{center}
\درج‌تصویر[پهنا=0.8\پهنای‌متن]{src/Theorical/tree.png}
\end{center}

%اگر مسئله چند قسمت داشت به این صورت عمل کنید:
%\شروع{زیرمسائل}
%\زیرمسئله[0] % مسئلۀ ۱آ
%\زیرمسئله % مسئلۀ ۱ب
%\زیرمسئله % مسئلۀ ۱پ
%\زیرمسئله % مسئلۀ ۱ت
%\پایان{زیرمسائل}

\solution{
	پاسخ نمونه.
}




%%%%%%%%%%%%%%%%%%%%%%%%%%%%%%%%%%%%%%%%%%%%%%%%%%%%%
% برخی دستورات مفید برای تمرین‌ها. %
%%%%%%%%%%%%%%%%%%%%%%%%%%%%%%%%%%%%%%%%%%%%%%%%%%%%%

% فرمول‌های نمونه:
%$f(x) = \Theta(x)$
%$T(x, y) \leq \log(x) + 2^y + \binom{2n}{n}$
% {\tt code\_function}


% لیست موارد بدون شماره‌گذاری:
%\شروع{فقرات}
%    \فقره مورد اول 
%        \شروع{فقرات}
%            \فقره مورد اول داخلی 
%                \شروع{فقرات}
%                    \فقره اولین مورد خیلی داخلی 
%                    \فقره دومین مورد خیلی داخلی 
%                \پایان{فقرات}
%            \فقره مورد دوم داخلی 
%        \پایان{فقرات}
%    \فقره مورد دوم 
%\پایان{فقرات}

% لیست موارد شماره‌گذاری شده:
%\شروع{شمارش}
%    \فقره یک مورد 
%    \فقره یک مورد دیگر 
%    \فقره باز هم یک مورد دیگر
%\پایان{شمارش}

% برای نوشتن یک فرمول چند خطی که یک شماره بخوره. برای حالت‌های دیگر به سایت زیر مراجعه کنید:
% https://www.overleaf.com/learn/latex/Aligning_equations_with_amsmath
%\begin{align} 
%    \begin{split}
%        (x+y)^3 &= (x+y)^2(x+y) \\
%                &= (x^2+2xy+y^2)(x+y) \\
%                &= (x^3+2x^2y+xy^2) + (x^2y+2xy^2+y^3) \\
%                &= x^3+3x^2y+3xy^2+y^3
%    \end{split}                                 
%\end{align}

% نوشتن ماتریس و آرایه دوبعدی
%\begin{align}
%    A = 
%    \begin{bmatrix}
%        A_{11} & A_{21} \\
%        A_{21} & A_{22}
%    \end{bmatrix}
%\end{align}

% اضافه کردن تصویر با شرح. برای اضافه کردن تصویر بدون شرح، فقط دستور درج‌تصویر را استفاده کنید.
%\شروع{شکل}[h!]
%\تنظیم‌ازوسط
%  \درج‌تصویر[پهنا=0.5\پهنای‌متن]{img.jpg}
%\شرح{شرح تصویر}
%\برچسب{تصویر:۱}
%\پایان{شکل}
%\پایان{زیرمسائل}

%کد زیر را ببینید.
%%\lstinputlisting{source_filename.py}
%\begin{latin}
%\begin{lstlisting}
%# Import the modules
%import sys
%import random
%
%ans = True
%
%while ans:
%question = raw_input("Ask the magic 8 ball a question: (press enter to quit) ")
%
%answers = random.randint(1,8)
%
%if question == "":
%sys.exit()
%
%elif answers == 1:
%print "It is certain"
%
%elif answers == 2:
%print "Outlook good"
%
%elif answers == 3:
%print "You may rely on it"
%
%elif answers == 4:
%print "Ask again later"
%
%elif answers == 5:
%print "Concentrate and ask again"
%
%elif answers == 6:
%print "Reply hazy, try again"
%
%elif answers == 7:
%print "My reply is no"
%
%elif answers == 8:
%print "My sources say no"
%\end{lstlisting}
%\end{latin}